\section{Introduction}
\label{sec:introduction}

LLM agents are increasingly deployed for complex, long-horizon software engineering tasks---resolving GitHub issues, debugging production systems, and iterating on multi-file codebases across dozens of editing rounds \citep{jimenez2024swebench,yang2024sweagent,wang2025openhands}.
As these deployments scale, a critical failure mode has emerged that is qualitatively different from crashes or syntax errors: the agent \emph{stagnates}---continuing to produce actions that appear purposeful while making no meaningful progress toward the task goal.
\citet{orlanski2026slopcodebench} document this degradation on SlopCodeBench, recording verbosity increases of 77\% and quality drops of 75.5\% in long-horizon coding trajectories, with no agent fully solving any benchmark problem end-to-end.
Crucially, the dominant failure pattern is not simple repetition but what we term \emph{semantic stagnation}: the agent issues varied commands, edits different files, and generates plausible-looking outputs, yet its behavior is functionally stuck.
This ``looks busy but going in circles'' pattern defeats existing monitoring tools, because surface-level action diversity masks the underlying lack of progress.
We observe this phenomenon at scale: across trajectories from two architecturally distinct agents---SWE-Agent \citep{yang2024sweagent} and OpenHands \citep{wang2025openhands}---on SWE-bench tasks \citep{jimenez2024swebench}, 28.5\% and 20.9\% of sliding windows exhibit stagnation under our silver labeling scheme, respectively, confirming that the failure mode is both prevalent and highly agent-dependent.

Existing approaches address aspects of this problem but fall short in three ways.
First, \emph{n-gram cycle detectors} \citep{george2025cycledetection,pipis2025reasoningloop} flag exact action-sequence repetition and match richer structural features on in-domain gold (${\sim}$0.84 macro-F1), but collapse on cross-agent transfer ($-$43.7\,pp; \S\ref{sec:gold_set})---they lack the richer structural vocabulary (fuzzy match ratios, edit distances, file operation counts) needed to generalize across agents.
Second, \emph{step-level scoring} via process reward models \citep{xi2025agentprm} evaluates individual actions in isolation, discarding the window-level temporal dynamics that distinguish transient setbacks from sustained stagnation---a distinction that requires reasoning over sequences of consecutive actions \citep{chen2025monitoring}.
A sequence of failed commands followed by a strategy switch, for instance, carries different diagnostic information than the same commands evaluated independently.
Third, \emph{LLM self-evaluation} \citep{shinn2023reflexion} achieves strong in-domain accuracy (0.910 gold macro-F1 with few-shot-CoT) but is costly (\$40+ per evaluation) and brittle across agents: the best LLM reaches only 0.491 on cross-agent transfer (\S\ref{sec:cross_agent}), ruling out deployment as a general-purpose monitor.
More fundamentally, none of these approaches asks the prior question: \emph{what signals actually carry stagnation information?}
No systematic study has investigated whether detecting subtle agent failure requires sophisticated semantic understanding or whether simpler structural patterns suffice.

We address this question through a systematic empirical study (Figure~\ref{fig:pipeline}).
Using a weak supervision framework \citep{ratner2016dataprogramming,ratner2019snorkel}, we construct silver-labeled datasets from over 205K windows spanning SWE-Agent (54,809 windows, 28.5\% silver-labeled stagnation) and OpenHands (150,652 windows, 20.9\% stagnation).
Three labeling functions provide complementary noisy votes: command repetition ($\lambda_1$) targets repetitive behavior, file diff delta ($\lambda_2$) targets output productivity, and outcome backpropagation ($\lambda_3$) propagates task-level outcomes backward through the trajectory.
Majority vote aggregation produces window-level binary labels with confidence scores.
Our window-level formulation ($k{=}5$ consecutive actions, stride 1) captures short-range temporal patterns---such as repeated failures followed by strategy changes---that step-level approaches cannot express, while providing the granularity to pinpoint \emph{when} stagnation begins and ends within a trajectory, enabling targeted interventions.
Against these silver labels, we systematically compare hand-crafted structural features, a neural embedding probe, and an LLM-as-judge baseline.

Our study yields three principal findings.

\textbf{Finding 1: Structural features uniquely enable zero-data cross-agent transfer.}
On gold, the best LLM judge achieves the highest non-circular accuracy (0.910 vs.\ structural GB 0.839), but LLM judges collapse on cross-agent transfer (0.491 vs.\ 0.967, $-$47.6\,pp).
Cycle detectors match structural GB on in-domain gold (${\sim}$0.84) but likewise collapse cross-agent ($-$43.7\,pp).
In zero-data transfer, structural features outperform the probe by 39.4\,pp; with fine-tuning, the probe surpasses structural features (0.976 vs.\ 0.967).

\textbf{Finding 2: Label quality, not LF selection, is the dominant transfer bottleneck.}
LFs exhibit temporal dynamics that reverse across agents ($\lambda_1$ leads by $-1.87$ steps on SWE-Agent, lags by $+3.32$ on OpenHands).
Fixing OH label version bias (stag rate 6.56\%$\to$21.4\%) improves 5\% fine-tuning F1 from 0.570 to 0.850; $\lambda_1$'s core value is coverage ($-29.4$\,pp when removed), not precision.

\textbf{Finding 3: No label model improves over majority vote.}
Five label models---three temporal and two generative \citep{safranchik2020linkedhmm,fu2020flyingsquid,ratner2019snorkelmetal}---all fail to surpass majority vote on downstream detection.
The bottleneck is LF quantity: three binary LFs produce too little inter-LF conflict for any method to exploit.

Our contributions are:
\begin{enumerate}
    \item \textbf{Weak supervision methodology for agent monitoring.} We design a silver-labeling pipeline (3 LFs + majority vote) over 205K+ windows from two coding agents, enabling systematic comparison of detection methods without manual annotation.
    \item \textbf{Label quality as transfer bottleneck.} We show that labeling functions exhibit agent-specific temporal reversals and that fixing label version bias (stag rate 6.56\%$\to$21.4\% on the embedding subset) is the dominant factor for cross-coding-agent transfer, informing future weak-supervision design on sequential agent data.
    \item \textbf{Label model negative finding.} We show that five label models---three temporal and two generative---fail to improve over majority vote for downstream detection. The bottleneck is LF quantity, not aggregation method: the 3-LF regime produces insufficient inter-LF conflict for any approach to add value.
    \item \textbf{Deployment-regime trade-off.} On gold, the best LLM judge achieves the highest non-circular accuracy (0.910 vs.\ structural GB 0.839), but LLM judges collapse on cross-agent transfer (0.491 vs.\ 0.967, $-$47.6\,pp) at 8{,}000$\times$ higher cost. Structural features uniquely enable zero-data cross-agent transfer, outperforming the probe by 39.4\,pp and cycle detectors by 43.7\,pp (no GPU, ${<}$5\,ms per window). With fine-tuning, the probe surpasses structural features (0.976 vs.\ 0.967), delineating when each approach is warranted.
\end{enumerate}